% Dati, stile e comandi della tesi, in un solo posto.

% Dati del documento e del frontespizio.
\newcommand{\TitoloTesi}{ChatCart: agente conversazionale con interfaccia generativa}
\newcommand{\AutoreTesi}{Edoardo De Piccoli}
\newcommand{\Matricola}{2101055}
\newcommand{\Relatore}{Prof.\ Tullio Vardanega}
\newcommand{\AnnoAccademico}{2025-2026}

\title{\TitoloTesi}
\author{\AutoreTesi}
\date{}

% Pagina, caratteri e spaziatura.
\usepackage[margin=3cm]{geometry}
\usepackage{fontspec}
\usepackage{lmodern}
\usepackage[italian]{babel}
\usepackage{microtype}
\usepackage{setspace}
\setstretch{1.25}
\usepackage{parskip}
\usepackage{enumitem}
% Elenchi compatti e indivisibili, uniti alla frase che li introduce: se non entrano nella pagina
% passano interi alla successiva.
\makeatletter
\setlist[itemize,enumerate]{nosep, beginpenalty=10000, before={\interlinepenalty=10000\relax\@itempenalty=10000\relax}}
\makeatother
\setlength{\emergencystretch}{3em}
% Nessuna riga isolata di un paragrafo in cima o in fondo alla pagina.
% Si impostano a inizio documento perché un pacchetto le riporta a 3000.
\AtBeginDocument{\widowpenalty=10000 \clubpenalty=10000 \displaywidowpenalty=10000}
\setcounter{secnumdepth}{2}
\setcounter{tocdepth}{2}

% Note a piè di pagina sempre in fondo alla pagina, anche quando il testo non la riempie.
\usepackage[bottom]{footmisc}

% Titolo dei capitoli su una riga, con il numero davanti e senza la parola «Capitolo».
\usepackage{titlesec}
\titleformat{\chapter}[hang]{\normalfont\huge\bfseries}{\thechapter}{0.8em}{}
\titlespacing*{\chapter}{0pt}{70pt}{40pt}

\usepackage{graphicx}
\usepackage[table,dvipsnames]{xcolor}

% Le figure sono file finiti in assets/, inclusi con \includegraphics alla loro larghezza.
\usepackage{caption}
\captionsetup{font=small}
% Una figura alta fino a tre quarti della pagina resta nel testo invece di finire su una pagina
% di soli inserti.
\renewcommand{\floatpagefraction}{0.75}
\renewcommand{\topfraction}{0.85}
\renewcommand{\bottomfraction}{0.85}
\renewcommand{\textfraction}{0.1}
% Gli inserti restano nella sezione o sottosezione che li introduce.
\usepackage[section]{placeins}
\AddToHook{cmd/subsection/before}{\FloatBarrier}
\usepackage{float}
\floatstyle{plain}
\newfloat{codelisting}{htbp}{lol}[chapter]
\floatname{codelisting}{Codice sorgente}
\usepackage{listings}
\lstset{
	basicstyle=\small\ttfamily,
	breaklines=true,
	columns=fullflexible,
	keepspaces=true,
	showstringspaces=false
}

% Tabelle: stile comune e larghezza complessiva pari al testo.
% Una tabella non si spezza mai tra due pagine: se non entra nello spazio rimasto,
% passa intera alla pagina successiva. Didascalia e identificativo sono argomenti
% dell'ambiente, così la didascalia compare sempre sotto la tabella.
% I margini standard delle colonne allineano sfondi e righe senza patch.
\usepackage{booktabs,array}
\usepackage{longtable}
\usepackage{needspace}
\definecolor{TabellaIntestazione}{HTML}{E2E5E9}
\definecolor{TabellaRiga}{HTML}{F3F4F6}
\newcolumntype{P}[1]{>{\raggedright\arraybackslash}p{\dimexpr#1\linewidth-2\tabcolsep\relax}}
\newenvironment{tabellatesi}[3]{%
	\def\TabellaDidascalia{#2}%
	\def\TabellaEtichetta{#3}%
	\begin{table}[htbp]%
	\centering
	\setstretch{1}\fontsize{10.5}{13}\selectfont
	\setlength{\tabcolsep}{8pt}%
	\renewcommand{\arraystretch}{1.2}%
	\setlength{\extrarowheight}{3pt}%
	\setlength{\aboverulesep}{0pt}%
	\setlength{\belowrulesep}{0pt}%
	\arrayrulecolor{black!25}%
	\captionsetup{position=below,skip=6pt}%
	\rowcolors{2}{white}{TabellaRiga}%
	\begin{tabular}{#1}%
}{%
	\end{tabular}%
	\caption{\TabellaDidascalia}\label{\TabellaEtichetta}%
	\end{table}%
}

% Tabelle del capitolo 3: righe intere, intestazione ripetuta quando cambia pagina.
% Argomenti: colonne, intestazione, didascalia, etichetta.
\newenvironment{tabellalunga}[4]{%
	\par\begingroup
	\def\TabellaDidascalia{#3}%
	\def\TabellaEtichetta{#4}%
	\setstretch{1}\fontsize{10.5}{13}\selectfont
	\setlength{\tabcolsep}{8pt}%
	\renewcommand{\arraystretch}{1.2}%
	\setlength{\extrarowheight}{3pt}%
	\setlength{\aboverulesep}{0pt}%
	\setlength{\belowrulesep}{0pt}%
	\setlength{\LTcapwidth}{\linewidth}%
	\arrayrulecolor{black!25}%
	\rowcolors{2}{white}{TabellaRiga}%
	\begin{longtable}{#1}%
	\toprule
	\rowcolor{TabellaIntestazione} #2 \\
	\midrule
	\endfirsthead
	\toprule
	\rowcolor{TabellaIntestazione} #2 \\
	\midrule
	\endhead
	\endfoot
	\endlastfoot
}{%
	\bottomrule
	\rowcolor{white}%
	\caption{\TabellaDidascalia}\label{\TabellaEtichetta}\\
	\end{longtable}%
	\endgroup
}

% Il capitolo 3 consente interruzioni tra i punti degli elenchi. Il gruppo locale
% conserva le impostazioni dei capitoli già approvati e delle parti successive.
\newcommand{\iniziocapitoloterzo}{%
	\begingroup
	\setlist[itemize,enumerate]{nosep,beginpenalty=10000,midpenalty=0,
		before={\interlinepenalty=0\relax}}%
}
\newcommand{\finecapitoloterzo}{\clearpage\endgroup}

\usepackage[unicode]{hyperref}
\usepackage{bookmark}
\usepackage{xurl}
\urlstyle{same}
\hypersetup{
	pdftitle={\TitoloTesi},
	pdfauthor={\AutoreTesi},
	pdflang={it-IT},
	colorlinks=true,
	linkcolor=Maroon,
	citecolor=Blue,
	urlcolor=Blue
}

% Comandi per rimandi e glossario.
% Riferimenti: il nome e il numero sono un unico collegamento cliccabile.
\newcommand{\figref}[1]{\hyperref[#1]{figura~\ref*{#1}}}
\newcommand{\tabref}[1]{\hyperref[#1]{tabella~\ref*{#1}}}
\newcommand{\secref}[1]{\hyperref[#1]{sezione~\ref*{#1}}}
\newcommand{\codref}[1]{\hyperref[#1]{codice~\ref*{#1}}}
\newcommand{\listoflistings}{\listof{codelisting}{Elenco dei codici sorgente}}

% Capitoli finali senza numero, presenti nell'indice e nei segnalibri PDF.
\newcommand{\capitolononnumerato}[1]{%
	\chapter*{#1}%
	\phantomsection
	\addcontentsline{toc}{chapter}{#1}%
}

% Il pedice G si inserisce esplicitamente solo alla prima occorrenza.
% Termine e pedice fanno parte dello stesso collegamento.
\newcommand{\gloss}[2]{\hyperref[glossario-#1]{#2\textsubscript{G}}}
\newcommand{\rimandoglossario}[2]{\hyperref[glossario-#1]{#2}}
\newcommand{\voceglossario}[3]{%
	\item[\protect\phantomsection\label{glossario-#1}#2] #3%
}
